\documentclass[11pt]{article}
\usepackage{ochamath}
\subjectcolor{7A4BA8}
\setsheet{複素関数}{オイラーの公式とフェーザ　演習 1}{https://math.ochanote.com/complex/euler-phasor/}
\namelinetrue
\begin{document}
\makesheethead{問題}
{\small 虚数単位は $j$（$j^2=-1$）で表す。}

\problem[極形式]
次の複素数を $re^{j\theta}$（$r>0,\ -\pi<\theta\le\pi$）の形で表せ。(3) は値を求めよ。
\[ \text{(1)}\ 1+j \qquad \text{(2)}\ {-\sqrt3} + j \qquad \text{(3)}\ (1+j)^8 \]
\workspace{45mm}

\problem[オイラーの公式の応用]
\begin{enumerate}[label=(\arabic*)]
\item $\cos\theta = \dfrac{e^{j\theta} + e^{-j\theta}}{2}$ を示せ。
\item (1) を使って $\cos^2\theta = \dfrac{1 + \cos 2\theta}{2}$ を導け。
\end{enumerate}
\workspace{50mm}

\problem[RC直列回路]
$R = 10\ \Omega$、$\dfrac{1}{\omega C} = 10\ \Omega$ の直列回路に $v(t) = 20\cos\omega t\ \mathrm{V}$ を加えた。フェーザを用いて電流 $i(t)$ を求めよ。
\workspace{55mm}

\problem[直列共振]
$R$、$L = 10\ \mathrm{mH}$、$C = 1\ \mu\mathrm{F}$ の直列回路について、
\begin{enumerate}[label=(\arabic*)]
\item インピーダンス $Z$ を $\omega$ の式で表せ。
\item $Z$ が実数になる角周波数 $\omega_0$ と、そのときの周波数 $f_0$ を求めよ。
\end{enumerate}
\workspace{40mm}
\end{document}
